\section*{Sets and Indices}

Let
\[
C = \{\text{corn},\text{wheat},\text{soybeans},\text{sorghum}\}
\]
denote the set of crops. No additional indexing sets are used in the implemented model.

\section*{Parameters}

For each crop \(c \in C\), let \(p_c\) denote profit per acre:
\[
p_{\text{corn}},\; p_{\text{wheat}},\; p_{\text{soybeans}},\; p_{\text{sorghum}}.
\]
In code these are
\texttt{profit\_corn}, \texttt{profit\_wheat}, \texttt{profit\_soybeans}, and \texttt{profit\_sorghum}.

Let
\[
A
\]
(code: \texttt{total\_area}) be the total available farm area.

Let
\[
r^{cw}
\]
(code: \texttt{corn\_wheat\_ratio}) be the minimum corn-to-wheat acreage ratio.

Let
\[
r^{ss}
\]
(code: \texttt{soybeans\_sorghum\_ratio}) be the minimum soybeans-to-sorghum acreage ratio.

Let
\[
r^{ws}
\]
(code: \texttt{wheat\_sorghum\_ratio}) be the wheat-to-sorghum proportionality factor.

\section*{Decision Variables}

Let
\[
x_{\text{corn}} \ge 0
\]
(code: \texttt{corn}) be the acres planted in corn.

Let
\[
x_{\text{wheat}} \ge 0
\]
(code: \texttt{wheat}) be the acres planted in wheat.

Let
\[
x_{\text{soybeans}} \ge 0
\]
(code: \texttt{soybeans}) be the acres planted in soybeans.

Let
\[
x_{\text{sorghum}} \ge 0
\]
(code: \texttt{sorghum}) be the acres planted in sorghum.

\section*{Objective}

Maximize total profit:
\[
\max \;
p_{\text{corn}} x_{\text{corn}}
+ p_{\text{wheat}} x_{\text{wheat}}
+ p_{\text{soybeans}} x_{\text{soybeans}}
+ p_{\text{sorghum}} x_{\text{sorghum}}.
\]

\section*{Constraints}

Total land availability:
\[
x_{\text{corn}} + x_{\text{wheat}} + x_{\text{soybeans}} + x_{\text{sorghum}} \le A.
\]

Corn-to-wheat ratio:
\[
x_{\text{corn}} \ge r^{cw} x_{\text{wheat}}.
\]

Soybeans-to-sorghum ratio:
\[
x_{\text{soybeans}} \ge r^{ss} x_{\text{sorghum}}.
\]

Wheat-to-sorghum proportionality:
\[
x_{\text{wheat}} = r^{ws} x_{\text{sorghum}}.
\]

Nonnegativity:
\[
x_{\text{corn}},x_{\text{wheat}},x_{\text{soybeans}},x_{\text{sorghum}} \ge 0.
\]

\section*{Notes on Implementation}

The implemented model is a continuous linear program with four nonnegative decision variables and one equality constraint. It is solved in \texttt{current\_heuristic.py} by constructing an \texttt{LpProblem} with maximization sense and passing it to \texttt{GUROBI\_CMD(msg=0)} through the PuLP-compatible interface.

A notable implementation detail is that the wheat--sorghum requirement is modeled as an equality,
\[
x_{\text{wheat}} = r^{ws} x_{\text{sorghum}},
\]
not as an inequality. This follows the source code exactly and should not be replaced by a weaker or stronger alternative when reproducing the model.

All acreage variables are continuous; there are no integrality restrictions.
